\section{Cấu hình GitHub}

\subsection{Fork Repository}
Nhóm đã fork repository gốc \texttt{nashtech-garage/yas} về tài khoản \texttt{tzin1401/yas}. Tất cả thành viên được thêm làm collaborator với quyền write.

% [Ảnh: Screenshot GitHub repo tzin1401/yas]

\subsection{Branch Protection Rules (GitHub Rulesets)}
Để đảm bảo chất lượng code, nhóm đã cấu hình GitHub Rulesets cho branch \texttt{main} với các quy tắc:

\begin{enumerate}[noitemsep]
    \item \textbf{Không cho phép push trực tiếp}: Mọi thay đổi phải thông qua Pull Request
    \item \textbf{Yêu cầu tối thiểu 2 reviewer approve}: Đảm bảo code review đầy đủ
    \item \textbf{Required status check}: CI pipeline phải pass (\texttt{continuous-integration/jenkins/pr-head})
    \item \textbf{Không cho phép force push}: Bảo vệ lịch sử commit
\end{enumerate}

% [Ảnh: Screenshot GitHub Rulesets configuration]
% [Ảnh: Screenshot hiển thị "Review required" và "Checks" trên PR]

\subsection{Webhook GitHub → Jenkins}
Cấu hình webhook để GitHub tự động thông báo Jenkins khi có push hoặc PR:
\begin{itemize}[noitemsep]
    \item \textbf{Payload URL}: \texttt{http://3.27.92.213:8080/github-webhook/}
    \item \textbf{Content type}: \texttt{application/json}
    \item \textbf{Events}: Push events, Pull request events
\end{itemize}

% [Ảnh: Screenshot GitHub Webhook configuration]
% [Ảnh: Screenshot Webhook delivery history (Recent Deliveries)]

\subsection{Quy trình làm việc với Git}
Nhóm tuân theo quy trình Git Flow đơn giản:
\begin{enumerate}[noitemsep]
    \item Tạo branch mới từ \texttt{main} (ví dụ: \texttt{feature/media-tests}, \texttt{ci/setup-jenkins-pipeline})
    \item Phát triển và commit code trên branch
    \item Push branch lên GitHub
    \item Tạo Pull Request vào \texttt{main}
    \item Jenkins tự động chạy CI pipeline
    \item 2 thành viên khác review và approve
    \item Merge PR khi CI pass và đủ approvals
\end{enumerate}
